\lecture{7}
\title{Statics and Grasping}

\begin{document}

\subsection{Forces: Spatial Force}

In this lecture, we'll answer the following question:
\begin{quote}
    \textbf{Question.}
    Given an object I want to grasp,
    where do I put my fingers,
    and will the object stay in my hand?
\end{quote}
This question involves not just geometry,
but the relevant \textbf{forces} being applied, too.

\textbf{Definition.}
We can represent \textbf{spatial force} via
a vector in $\mathbb{R}^6$ given by:
\[
F^{B_p}_C = \begin{bmatrix}
\tau_C^{B_p} \\ f_C^{B_p}
\end{bmatrix} ~ \text{ where } 
\tau^{B_p}_C \in \mathbb{R}^3 \text{ is torque, and }
f^{B_p}_C \in \mathbb{R}^3 \text{ is translational force.}
\]
where $B_p$ is the point on the body $B$ where the force is applied.
(As with ${}^AV^B_F$, angular stuff goes above translational stuff.)

\begin{theorem}{Spatial Force Properties}
    We can re-express spatial force in different frames like so:
    \[
    F^{B_p}_B = {}^BR^C F^{B_p}_C
    := \begin{bmatrix}
    {}^BR^C & 0 \\ 0 & {}^BR^C
    \end{bmatrix} F^{B_p}_C.
    \]
    Furthermore, if a force $i$ acts on body $B$
    at position $B_{c_i}$, then we get:
    \[
    f^B_{i, B} = f^{B_{c_i}}_{i, B}
    ~~~ \text{ and } ~~~
    \tau^B_{i, B} = \tau^{B_{c_i}}_{i, B}
    + {}^Bp^{B_{c_i}} \times f^{B_{c_i}}_{i, B}.
    \]
    There's nothing new or complicated here beyond AP Physics I, say.
\end{theorem}

\begin{remark}
    It is sometimes convenient to express the cross-product $d \times f$
    as a matrix multiplication $S(d) f$, given by:
    \[
    d \times f = \underbrace{\begin{bmatrix}
    0 & -d_z & d_y \\ d_z & 0 & -d_x \\ -d_y & d_x & 0
    \end{bmatrix}}_{S(d)} f.
    \]
\end{remark}

It then follows that a body is in \textbf{static equilibrium}
if the sum of the spatial forces acted upon it is zero.

\begin{remark}
    The reason why we formalize everything using spatial forces
    is it naturally gets us to express everything---particularly
    torques---in the same coordinate frame $B$.
\end{remark}

For example, we can model two hands grasping on
opposite faces of a box like so.
\begin{center}
    \frame{\includegraphics[width=0.4\textwidth]{lecture-07/grasp-forces.png}}
\end{center}
To achieve static equilibrium, note that
both the translational forces \emph{and}
torques about $B$ must sum to zero.

\begin{remark}
    It suffices to only think about statics (i.e., no $F = ma$)
    because at manipulation speeds,
    accelerations $a$ are far below $g$.
\end{remark}

\subsection{Modeling Contacts}

We only model contacts at a \emph{point on a plane}
(as opposed to, say, a point along an edge).

\begin{remark}
    If our contact is not a point---for example,
    the contact between a box and the table it lies on---then
    we can approximate that contact as
    the sum of point contacts---for example,
    the sum of contacts at the four vertices
    of the bottom face of the box.
\end{remark}

There are three kinds of point-on-plane contacts.
\begin{itemize}
    \item \textbf{Frictionless.}
    Forces can only be applied perpendicular to the plane.
    \item \textbf{Point w/ Friction.}
    Forces can have normal and tangential components,
    but their ratio must be below $\mu$.
    \begin{center}
        \frame{\includegraphics[width=0.2\textwidth]{lecture-07/force-cone.png}}
    \end{center}
    Equivalently, forces can only be applied inside the friction cone;
    analytically, $\sqrt{(f^{c_i}_{C_x})^2 + (f^{c_i}_{C_y})^2}
    \leq \mu f^{c_i}_{C_z}$.
    \begin{itemize}
        \item For the sake of \emph{linear algebra},
        we oftentimes approximate (discretize) the friction cone as a pyramid.
    \end{itemize}
    \item \textbf{Soft Finger.}
    In addition to forces within the friction cone,
    you can also apply torques $\tau_{\mathrm{normal}}$ locally
    if $|\tau_{\mathrm{normal}}| \leq \gamma f^{c_i}_{C_z}$.
\end{itemize}
There are some obvious limitations to these models---friction
is more complicated than just ``$\mu$'',
and the patch of contact is not just a single point, and so on\dots

\subsection{Grasping: Is It Possible?}

First, some notation that we use for
mapping multiple applied forces with different contact frames
into the same body frame.

\textbf{Definition.}
Say the orientation of the $i^{\text{th}}$ contact frame is given by
$Q^{(i)} := {}^BR^{C_i} \in \mathbb{R}^{3 \times 3}$,
meaning the $i^{\text{th}}$ contact frame
applies a force $f_i \in \mathbb{R}^3$ on body $B$ at position $C_i$.
Then the force \emph{in the frame of $B$} can be expressed as:
\[
F = \sum_{i = 1}^k G_i f_i = G \begin{bmatrix}
f_1 \\ \vdots \\ f_k
\end{bmatrix}
~~~~~ \text{ where } ~~~~~
G_i := \begin{bmatrix}
S^{(i)}Q^{(i)} \\ Q^{(i)}
\end{bmatrix} \in \mathbb{R}^{6 \times 3}
~ \text{ and } ~
G := \begin{bmatrix} G_1 & \cdots & G_k \end{bmatrix}
~ \text{ is the } \textbf{grasp map.}
\]
The idea is that each matrix $G_i$
is the appropriate map from forces $f_i$ expressed in the contact frame
to forces expressed in the body's frame.
Here $S^{(i)} := S({}^Bp^{C_i}) \in \mathbb{R}^{3 \times 3}$
is the ``cross-product matrix'' mentioned in a remark earlier.

The goal of grasping is to answer the following question:
\begin{quote}
    \textbf{Question.}
    Given external forces $F^{\mathrm{ext}}$,
    is there an \emph{admissible} grasping force $f$
    such that $Gf + F^{\mathrm{ext}} = 0$ holds?

    (A force $f$ is \emph{admissible} if it lives within the friction cone.)
\end{quote}

\begin{remark}
    The grasp map can have a nontrivial null space;
    two grasping methods can both yield equilibrium.
    \begin{center}
        \frame{\includegraphics[width=0.4\textwidth]{lecture-07/grasp-null-space.png}}
    \end{center}
    In the diagram above, for example,
    the hard squeeze is just the light squeeze plus a
    ``left-right'' squeeze $f_0$ for which $Gf_0 = 0$.
\end{remark}

\textbf{Definition.}
The \textbf{grasp wrench space} is the set:
\[
    \mathcal{W} :=
    \left\{ F \mid F = \sum_{i = 1}^k G_i f_i
    \text{ for some } f_i \in \mathcal{K}_i \right\},
\]
where $\mathcal{K}_i$ is the space of admissible grasping forces.
(Note that this is the \emph{Minkowski sum} of the sets $G_i \mathcal{K}_i$.)

To answer the question from before for \emph{every} $F^{\mathrm{ext}}$,
it suffices to simply compute $\mathcal{W}$
and determine whether $0 \in \mathcal{W}$
(not on its boundary!).
Since each $\mathcal{K}_i$ is convex,
we can estimate $\mathcal{W}$ easily by \dots
\begin{itemize}
    \item \dots approximating each cone $\mathcal{K}_i$
    as a finite-vertex convex polyhedron, then \dots
    \item \dots taking the Minkowski sum of the vertices,
    then the convex hull of the Minkowski sum.
\end{itemize}

\subsection{Closure}

\textbf{Definition.}
Preferably, we want grasps that can hold bodies in place.
More specifically, a grasp can have\dots
\begin{itemize}
    \item \textbf{Form Closure.}
    The contacts are geometrically positioned
    such that the body cannot move without penetrating a contact.
    \item \textbf{Force Closure.}
    The contact forces have a $\mathcal{W}$
    universal enough to resist any $F^{\mathrm{ext}}$
    (i.e., external force and torque).
\end{itemize}
Check with yourself that the following make sense;
dimensionality arguments should suffice.
\begin{center}
    \frame{\includegraphics[width=0.4\textwidth]{lecture-07/closure-table.png}}
\end{center}
But there's more to a grasp than whether it has force closure or not.

\textbf{Heuristic.} Given a grasp with force closure
and wrench space $\mathcal{W}$, consider:
\begin{itemize}
    \item The ``worst case'' quantity
    $\epsilon := \min_{F \in \partial\mathcal{W}} \|F\|$.
    \item The ``average'' quantity
    $V := \text{the volume of } \mathcal{W}$.
\end{itemize} 
We prefer grasps with larger $\epsilon$,
as the strength of the contact force required trends as $O(1/\epsilon)$.
Similarly, we prefer grasps with larger $V$.
By this metric, we can justify why,
among the two below grasps with force closure,
the grasp on the right is preferred.
\begin{center}
    \frame{\includegraphics[width=0.35\textwidth]{lecture-07/grasp-heurestics.png}}
\end{center}

\subsection{Heuristic Grasping}

For general bodies, what are some easy ways
we can design grasps that are guaranteed to have force closure?

\textbf{Definition.}
Given a grasp with two contact forces,
we say the grasp is \textbf{antipodal} if
the line connecting the contact points
lies within both contact forces' friction cones.
\begin{center}
    \frame{\includegraphics[width=0.4\textwidth]{lecture-07/antipodal.png}}
\end{center}
\begin{theorem}{Antipodal}
    All antipodal grasps in $\mathbb{R}^2$
    and all antipodal grasps in $\mathbb{R}^3$
    with \emph{soft finger contacts} have force closure.
\end{theorem}

This inspires the following heuristic
for generating grasps with force closure.

\begin{enumerate}
    \item Using simple robotic vision,
    determine a point cloud of the surface of the object.
    \item Use PCA to determine normal vectors
    perpendicular to the surface at each point in the point cloud.
    \item Pick a point $P$ at random,
    and follow its normal toward the inside of the body
    until it meets a point $P'$ on the surface.
    \item Check if a grasp at $P$ and $P'$ is valid
    (i.e., is antipodal and does not have intersection between
    the body and the gripper).
\end{enumerate}

This works pretty well in practice, too!

\begin{remark}
    In practice, you probably want to do more than
    just randomly pick among antipodal grasps;
    more heuristics are used.
\end{remark}

\subsection{Learned Grasping}

You can also train a CNN that
proposes grasps and learns to score
whether they are good (e.g., antipodal).
More generally, there are many families of
models for scoring grasps.
We won't get into the technical details.

\end{document}